% MUL poster starter — an institutional A0 poster layout without bundled marks.
\documentclass[final]{beamer}
\usepackage[size=a0,scale=1.12]{beamerposter}
\usepackage{fontspec}
\setmainfont{TeX Gyre Heros}
\usepackage{booktabs,tikz}
\usetikzlibrary{arrows.meta,positioning}
\definecolor{mulblack}{HTML}{202124}
\definecolor{mulorange}{HTML}{E87511}
\definecolor{mulgray}{HTML}{F3F4F5}
\setbeamertemplate{navigation symbols}{}
\setbeamercolor{normal text}{fg=mulblack,bg=white}
\setbeamercolor{block title}{fg=white,bg=mulblack}
\setbeamercolor{block body}{fg=mulblack,bg=mulgray}
\setbeamertemplate{block begin}{\begin{beamercolorbox}[wd=\linewidth,sep=0.65em]{block title}\Large\bfseries\insertblocktitle\end{beamercolorbox}\nointerlineskip\begin{beamercolorbox}[wd=\linewidth,sep=0.9em]{block body}}
\setbeamertemplate{block end}{\end{beamercolorbox}\vspace{0.9em}}
\title{Materials that remember the load history}
\author{Lea Novak · Materials and Interfaces Group · Montanuniversität-style poster starter}
\institute{Example laboratory · replace this line with your affiliation}
\begin{document}
\begin{frame}[t]
\begin{beamercolorbox}[wd=\paperwidth,sep=1.7em,leftskip=2em,rightskip=2em]{block title}
  {\fontsize{43}{50}\selectfont\bfseries\inserttitle}\par\vspace{0.25em}
  {\Large\insertauthor\hfill\insertinstitute}
\end{beamercolorbox}
\vspace{1.2em}
\begin{columns}[t,totalwidth=0.96\paperwidth]
\begin{column}{0.31\paperwidth}
\begin{block}{Question}
  Can a compact history descriptor predict fatigue damage when the load path
  contains long quiet intervals and short high-amplitude bursts?
\end{block}
\begin{block}{Experimental design}
  \begin{enumerate}\itemsep0.5em
    \item cycle three representative loading histories;
    \item measure stiffness after every 500 cycles;
    \item compare the descriptor with a full-history baseline.
  \end{enumerate}
\end{block}
\begin{block}{Contribution}
  The descriptor stores the turning points and the accumulated dwell time. It
  is small enough to inspect, yet retains the information that matters for
  the failure forecast.
\end{block}
\end{column}
\begin{column}{0.37\paperwidth}
\begin{block}{Method and signal}
\centering
\begin{tikzpicture}[x=1.25cm,y=1cm,>=Latex]
  \draw[->,thick] (0,0)--(9,0) node[right]{cycle};
  \draw[->,thick] (0,0)--(0,4.8) node[above]{stress};
  \draw[mulorange,very thick] plot[smooth] coordinates{(0,1.1)(1,1.3)(2,1.2)(3,1.6)(4,4.0)(5,1.5)(6,1.6)(7,3.3)(8,1.7)};
  \draw[mulblack,dashed,thick] plot[smooth] coordinates{(0,1.2)(1,1.3)(2,1.3)(3,1.5)(4,2.2)(5,1.6)(6,1.7)(7,2.1)(8,1.8)};
  \node[fill=white,draw=mulorange] at (6.6,4.2) {measured};
  \node[fill=white,draw=mulblack] at (6.7,3.55) {descriptor};
\end{tikzpicture}
\vspace{0.5em}
\begin{alertblock}{Headline result}
  The history descriptor predicts the next stiffness drop with 91\% accuracy,
  using 14\% of the storage required by the full signal.
\end{alertblock}
\end{block}
\begin{block}{Comparison}
\centering\begin{tabular}{lrrr}\toprule
method & storage & accuracy & lead time\\\midrule
full history & 100\% & 0.94 & 12 h\\
turning points & 14\% & 0.91 & 11 h\\
mean + variance & 4\% & 0.73 & 5 h\\\bottomrule
\end{tabular}
\end{block}
\end{column}
\begin{column}{0.25\paperwidth}
\begin{block}{What to remember}
  \begin{itemize}\itemsep0.55em
    \item Preserve the events that change the prediction.
    \item Keep the representation small enough to audit.
    \item Report storage and accuracy together.
  \end{itemize}
\end{block}
\begin{block}{Limitations}
  The example covers one alloy and one temperature range. The descriptor may
  need an additional state variable when the load history changes sign often.
\end{block}
\begin{block}{Next step}
  Test transfer to a second alloy, preregister the failure threshold, and make
  the history representation available with the measurement data.
\end{block}
\end{column}
\end{columns}
\vfill
\begin{beamercolorbox}[wd=\paperwidth,sep=0.8em,leftskip=2em,rightskip=2em]{block title}
  \small Code and data: example.org/project \hfill A0 landscape/portrait options can be changed above
\end{beamercolorbox}
\end{frame}
\end{document}
